% Agentic Sleep paper draft (LaTeX version)
% NeurIPS-ready preprint setup with article fallback for local compilation.
%
% Usage (preferred, NeurIPS style):
%   Place neurips_2024.sty in this directory, then compile:
%   pdflatex agentic_sleep_paper_draft.tex
%
% Usage (fallback article mode):
%   Compile directly without neurips_2024.sty.

\newif\ifneuripsstyle
\IfFileExists{neurips_2024.sty}{\neuripsstyletrue}{\neuripsstylefalse}

\ifneuripsstyle
  \documentclass{article}
  \usepackage[preprint]{neurips_2024}
\else
  \documentclass[11pt]{article}
  \usepackage[margin=1in]{geometry}
  \usepackage{setspace}
  \usepackage[numbers]{natbib}
  \onehalfspacing
\fi

\usepackage[T1]{fontenc}
\usepackage[utf8]{inputenc}
\usepackage{lmodern}
\usepackage{microtype}
\usepackage{amsmath,amssymb,amsfonts}
\usepackage{booktabs}
\usepackage{array}
\usepackage{hyperref}
\usepackage{xcolor}
\usepackage{enumitem}
\usepackage{url}
\usepackage{verbatim}

\hypersetup{
  colorlinks=true,
  linkcolor=blue!50!black,
  citecolor=blue!50!black,
  urlcolor=blue!50!black
}

\title{Agentic Sleep: Off-Cycle Memory Consolidation for Long-Horizon Autonomous Agents}

\author{
  Riyad Mehdiyev \\
  \texttt{[affiliation / lab]} \\
  \texttt{[email]}
  \And
  Collaborators \\
  \texttt{[to be filled]}
}

\date{}

\begin{document}

\maketitle

\begin{abstract}
Long-horizon autonomous agents degrade as interaction histories grow, even when large context windows are available. The primary failure mode is not only token budget exhaustion but also contextual noise: irrelevant observations, transient errors, and repeated reasoning traces that flatten action distributions and increase uncertainty. We present \textbf{Agentic Sleep}, an off-cycle memory consolidation framework that periodically converts raw interaction traces into compact heuristic rules and a lightweight semantic graph, then clears the active narrative buffer. The framework is implemented as a prototype with (i) entropy-guided sleep triggering, (ii) a three-tier memory hierarchy (short-term interaction, episodic, and long-term semantic memory), (iii) Bayesian signal/noise filtering for heuristic promotion and pruning, and (iv) a failure-seeded synthetic ``Dream'' phase for heuristic hardening.

We evaluate the current implementation in a controlled simulation of long-horizon logic-game-like tasks (Codenames-like and murder-mystery-like domains) and observe improved token efficiency and higher task performance relative to a raw-context control. These results are preliminary and use a simulated policy/environment, not a production LLM backend. We provide a paper-ready experimental protocol and result tables with placeholders for future LLM-based validation.
\end{abstract}

\section{Introduction}

Autonomous LLM agents are often evaluated on short-horizon tasks, yet many practical workloads (software engineering, multi-stage research, incident response, and strategic games) require coherent reasoning across tens to hundreds of turns. Contemporary agent designs frequently combine reasoning traces and tool use patterns (e.g., CoT/ReAct-style prompting) with iterative self-improvement loops \citep{wei2022cot, yao2023react, madaan2023selfrefine}. A common engineering response to long-horizon degradation is to increase context length and maintain more raw history. In practice, this introduces a different bottleneck: the model must repeatedly process stale, redundant, and low-value traces alongside task-critical evidence. The result is often semantic drift, delayed correction of invalid assumptions, and unstable action selection.

This paper studies a complementary approach: \textbf{structural abstraction over raw retention}. Instead of preserving the full trajectory in the active context, the agent periodically enters an off-cycle \emph{Sleep Phase} that consolidates recent interaction logs into high-density, reusable heuristics and a semantic memory graph. The agent then resumes operation using these abstractions plus a short recency window, rather than the full narrative history.

The central claim is that, for long-horizon tasks, a memory system that compresses \emph{what works} into reusable procedural rules can outperform a raw-context baseline in both token efficiency and task success.

\section{Problem Statement}

Let an autonomous agent operate over a sequence of turns $t = 1, \dots, T$. At each turn, it receives observations, produces reasoning traces, chooses an action, and obtains a task outcome. In raw-context architectures, the agent carries an increasingly long transcript into subsequent decisions. As the transcript grows, task-irrelevant tokens accumulate and the effective signal-to-noise ratio decreases.

We define the operational problem as:

\begin{itemize}[leftmargin=1.25em]
\item Maintain or improve decision quality over long horizons.
\item Reduce inference token cost relative to full-history control.
\item Prevent repeated integration of transient errors into long-term strategy.
\item Preserve historical lineage when knowledge changes (e.g., contradictory facts).
\end{itemize}

\section{Agentic Sleep Framework}

\subsection{Overview}

Agentic Sleep introduces a cyclical control loop:

\begin{enumerate}[leftmargin=1.25em]
\item \textbf{Waking Phase}: The agent performs the task using recency memory plus top-ranked heuristics.
\item \textbf{Trigger Check}: Sleep is triggered by elapsed turns and/or an entropy spike.
\item \textbf{Consolidation Phase}: Recent traces are converted into heuristic rules and semantic graph edges.
\item \textbf{Dream Phase}: Newly consolidated rules are stress-tested in synthetic counterfactual scenarios.
\item \textbf{Structural Reset}: The narrative buffer is cleared while compact rules and semantic summaries are retained.
\end{enumerate}

\subsection{Memory Hierarchy}

The prototype implements a three-tier memory hierarchy:

\begin{itemize}[leftmargin=1.25em]
\item \textbf{Tier 1: Short-Term Interaction Memory (STIM)} --- linear buffer of recent \texttt{EpisodeStep} objects used for recency-sensitive decisions.
\item \textbf{Tier 2: Mid-Term Episodic Memory (MTEM)} --- structured episodes formed by flushing STIM during sleep cycles.
\item \textbf{Tier 3: Long-Term Semantic Memory (LTSM)} --- a heuristic rule store plus a lightweight bi-temporal knowledge graph.
\end{itemize}

This design allows the agent to preserve recency, store audit-ready episodes, and maintain abstract strategic knowledge separately.

\subsection{Entropy-Guided Sleep Trigger}

The framework tracks Shannon entropy over the action probability distribution at each turn. Rising entropy is treated as a proxy for uncertainty or drift. Sleep can be triggered by:

\begin{itemize}[leftmargin=1.25em]
\item a fixed interval (e.g., every 10 turns), or
\item an entropy spike above threshold $\Delta H > \tau$, with a minimum gap to prevent thrashing.
\end{itemize}

This creates an adaptive reset mechanism: the agent consolidates when its decision distribution becomes unstable, not only on a rigid schedule.

In the current prototype, entropy is computed over action distributions rather than token distributions; the final LLM evaluation will replace this with next-token (or action-logit) entropy when provider outputs support it. The uncertainty metric follows Shannon entropy as a general measure of distributional uncertainty \citep{shannon1948mathematical}.

\section{Consolidation Mechanism}

\subsection{Input Representation}

Each interaction step is recorded as an \texttt{EpisodeStep} containing:

\begin{itemize}[leftmargin=1.25em]
\item turn index
\item domain label
\item compact thought trace (logging artifact)
\item selected action
\item observation/outcome string
\item success flag
\item entropy
\item semantic tags
\end{itemize}

STIM is flushed into an \texttt{Episode} during sleep. Episodes are retained in MTEM for traceability and future analysis.

\subsection{Heuristic Extraction}

The consolidation engine transforms episodic traces into procedural heuristics of the form:

\begin{quote}
When observing tag set $X$, prefer action $A$.
\end{quote}

In the current prototype, heuristics are stored as $(\text{action}, \text{tag})$ rules with Bayesian success statistics:

\begin{itemize}[leftmargin=1.25em]
\item \texttt{alpha}, \texttt{beta} initialize a Beta posterior
\item $p_{\text{success}} = \alpha / (\alpha + \beta)$
\item \texttt{support} tracks empirical evidence count
\item \texttt{repetitiveness} estimates low-information repetition in retained examples
\end{itemize}

This supports signal selection and pruning without requiring gradient updates.

\subsection{Signal vs.\ Noise Selection}

After procedure extraction, rules are categorized:

\begin{itemize}[leftmargin=1.25em]
\item \textbf{Promote (signal)} if posterior success exceeds a threshold and repetitiveness remains below a threshold.
\item \textbf{Prune (noise)} if support is sufficient but posterior success is consistently low.
\end{itemize}

This operationalizes the noise-signal dilemma in a lightweight Bayesian filter.

\subsection{Bi-Temporal Semantic Graph}

The long-term semantic memory also stores fact-like edges with timestamps:

\begin{itemize}[leftmargin=1.25em]
\item system-time creation/expiration
\item validity/invalidation timestamps
\item provenance episode ID
\end{itemize}

When a new edge contradicts an active edge with the same $(\text{subject}, \text{predicate})$, the previous edge is invalidated rather than deleted. This preserves lineage while maintaining an active current view of semantic memory.

\section{Dream Phase: Failure-Seeded Synthetic Hardening}

After consolidation, the Dream phase performs internal synthetic tests on top heuristics.

\subsection{Failure Seeding}

Dream uses recent failed steps (and optionally high-entropy steps) as seeds. Seed tags define the neighborhoods where the agent's current strategy appears brittle.

\subsection{Counterfactual Simulation}

For each supported top-ranked rule, the simulator samples counterfactual tags (preferably from failure seeds) and evaluates whether the rule still appears successful under nearby conditions. This produces synthetic support and synthetic success counts.

\subsection{Confidence Update}

Each rule's confidence is updated by mixing:

\begin{itemize}[leftmargin=1.25em]
\item empirical posterior success (from waking episodes), and
\item synthetic success rate (from Dream self-play)
\end{itemize}

This changes future rule influence without rewriting the raw episodic record.

\section{Prototype Implementation}

The current implementation is a runnable research prototype in Python (stdlib-only), designed to later accept a real LLM backend. The orchestration is LangGraph-inspired, using a cyclic state machine with conditional routing and memory-aware state transitions \citep{langgraph_docs}. Key modules:

\begin{itemize}[leftmargin=1.25em]
\item \texttt{agentic\_sleep/engine.py}: agent loop, scoring, entropy, triggers, sleep cycle
\item \texttt{agentic\_sleep/memory.py}: STIM/MTEM/LTSM hierarchy
\item \texttt{agentic\_sleep/consolidation.py}: heuristic extraction, graph updates, HCR computation
\item \texttt{agentic\_sleep/dream.py}: synthetic hardening
\item \texttt{agentic\_sleep/experiment.py}: control vs.\ sleep benchmark harness
\item \texttt{agentic\_sleep/langgraph\_adapter.py}: optional LangGraph-compatible graph skeleton
\end{itemize}

\subsection{Decision Policy (Current Prototype)}

The current agent policy is simulated (not an LLM) and combines:

\begin{itemize}[leftmargin=1.25em]
\item heuristic contributions from LTSM (sleep mode only),
\item recent episodic evidence from raw or STIM history,
\item a context-noise penalty that increases with effective context size,
\item a stochastic sampling step over normalized action scores.
\end{itemize}

This setup is intentionally simple: it isolates memory architecture effects before integrating external model variability.

\section{Related Work and Positioning}

Agentic Sleep is adjacent to, but distinct from, several lines of prior work. Reflexion-style agents use verbal feedback to improve subsequent attempts, storing textual reflections as episodic memory \citep{shinn2023reflexion}. Agentic Sleep differs by moving consolidation to an off-cycle phase that restructures memory into heuristics and semantic graph edges, and by explicitly pruning low-value procedures. MemGPT emphasizes paging and virtual memory management for raw information across memory tiers \citep{packer2023memgpt}; our framework instead prioritizes trajectory abstraction into procedural rules. Recursive summarization compresses long histories into compact narratives \citep{wang2023recursive_summary_memory}; Agentic Sleep replaces narrative summaries with condition-action heuristics to improve reuse across task instances.

At the prompting level, our framework is compatible with CoT/ReAct style traces and self-refinement loops \citep{wei2022cot, yao2023react, madaan2023selfrefine}. The contribution is therefore primarily architectural: an explicit off-cycle consolidation and hardening mechanism with measurable compression and efficiency metrics.

\section{Experimental Protocol (Current Prototype)}

\subsection{Task Environments}

We evaluate on synthetic long-horizon task families designed to mimic logic-heavy coordination and deduction:

\begin{itemize}[leftmargin=1.25em]
\item \textbf{Codenames-like domain}: clueing/guessing strategy abstractions and ambiguity management
\item \textbf{Murder-mystery-like (MMG) domain}: contradiction checking, timeline reconstruction, scene reinspection
\end{itemize}

Each turn is represented by a tag set and a hidden best action. Difficulty and distractor tags create ambiguity.

\subsection{Real Benchmark Tasks and Environments (Final Evaluation Plan)}

Beyond the synthetic domains used for prototype validation, the final study should evaluate Agentic Sleep on executable and contamination-aware benchmarks:

\begin{itemize}[leftmargin=1.25em]
\item \textbf{Software engineering issue resolution}: SWE-bench / SWE-bench Lite / SWE-bench Verified for repository-level bug fixing with execution-based validation \citep{jimenez2024swebench}.
\item \textbf{Code generation and repair micro-benchmarks}: HumanEval-style problems for fast iteration and ablations on memory/control decisions \citep{chen2021codex}.
\item \textbf{Interactive tool-using environments}: ReAct-style benchmark environments (e.g., ALFWorld/WebShop-style setups) to test long-horizon memory under observation-action loops \citep{yao2023react}.
\end{itemize}

In this repository, we provide a benchmark-file adapter (JSON/JSONL replay format) so real environment trajectories can be converted into turn-level evaluation instances and replayed under matched control vs.\ sleep conditions. This enables a clean transition from the current synthetic testbed to real agent traces without rewriting the core memory pipeline.

\subsection{Comparison Conditions}

\paragraph{Control (Raw Context).}
\begin{itemize}[leftmargin=1.25em]
\item Retains full raw trajectory for scoring
\item No consolidation
\item No sleep/dream phase
\end{itemize}

\paragraph{Sleep (Agentic Sleep).}
\begin{itemize}[leftmargin=1.25em]
\item Uses STIM + top heuristics + LTSM summaries
\item Triggers sleep by interval and entropy spike
\item Consolidates and runs Dream hardening
\item Clears active narrative buffer post-sleep
\end{itemize}

\subsection{Metrics}

Primary metrics:

\begin{itemize}[leftmargin=1.25em]
\item \textbf{Win Rate}
\item \textbf{Deductive Accuracy}
\item \textbf{Token Budget} (approximate inference token accounting)
\item \textbf{Heuristic Compression Ratio (HCR)}
\begin{equation}
\mathrm{HCR} = \frac{\mathrm{tokens}(\mathrm{raw\ traces})}{\mathrm{tokens}(\mathrm{consolidated\ heuristics + summaries})}
\end{equation}
\item \textbf{Consolidation Efficiency}
\begin{equation}
E_c = \frac{S_{\mathrm{sleep}} - S_{\mathrm{control}}}{B_{\mathrm{sleep}} / B_{\mathrm{control}}}
\end{equation}
where $S$ is success (here: win rate) and $B$ is token budget.
\end{itemize}

\section{Preliminary Results (Simulation Only; Non-LLM)}

The following run was produced by the current prototype:

\begin{verbatim}
python3 -m agentic_sleep.cli --games 20 --turns 50
\end{verbatim}

\subsection{Aggregate Results}

\begin{table}[t]
\centering
\caption{Aggregate prototype results (simulation only; non-LLM).}
\label{tab:aggregate-results}
\begin{tabular}{lrr}
\toprule
Metric & Control & Sleep \\
\midrule
Games & 40 & 40 \\
Win Rate & 0.325 & 0.475 \\
Deductive Accuracy & 0.551 & 0.601 \\
Avg Entropy & 0.673 & 0.246 \\
Token Budget & 136,429,271 & 28,372,395 \\
Sleep Cycles & 0 & 9 \\
Consolidations & 0 & 9 \\
Avg HCR & N/A & 7.49 \\
\bottomrule
\end{tabular}
\end{table}

Derived metric:
\begin{itemize}[leftmargin=1.25em]
\item \textbf{Consolidation Efficiency $E_c$}: \textbf{0.721}
\end{itemize}

\subsection{Domain Breakdown}

\begin{table}[t]
\centering
\caption{Domain-level prototype results (simulation only; non-LLM).}
\label{tab:domain-breakdown}
\begin{tabular}{lrrrr}
\toprule
Domain & Ctrl Win & Sleep Win & Ctrl Acc & Sleep Acc \\
\midrule
Codenames-like & 0.25 & 0.55 & 0.539 & 0.627 \\
MMG-like & 0.40 & 0.40 & 0.563 & 0.575 \\
\bottomrule
\end{tabular}
\end{table}

\subsection{Interpretation}

These preliminary results support the design hypothesis in the synthetic setting:

\begin{itemize}[leftmargin=1.25em]
\item Sleep improves overall performance while reducing token cost.
\item Entropy is lower in the sleep condition, consistent with reduced decision diffusion.
\item Gains are strongest in the Codenames-like domain, suggesting heuristic abstraction is especially effective when recurring semantic patterns map well to strategic actions.
\end{itemize}

However, these results \textbf{do not yet validate performance on real LLM agents}. They only indicate that the memory architecture is promising under a controlled simulated policy.

\section{Planned LLM-Based Evaluation (To Be Filled)}

This section is intentionally structured for later insertion of real-model results.

\subsection{LLM Integration Plan}

Replace the simulated policy with an LLM-backed action policy while preserving the same memory pipeline:

\begin{itemize}[leftmargin=1.25em]
\item Keep STIM/MTEM/LTSM unchanged
\item Replace the score-based action selector with:
  \begin{itemize}[leftmargin=1.25em]
  \item prompt construction from STIM + top rules + task state
  \item model inference
  \item action parsing + confidence estimation
  \end{itemize}
\item Use logprobs (if available) to compute entropy directly
\item Keep the same sleep trigger, consolidation, and Dream phases
\end{itemize}

\subsection{Proposed LLM Experiment Matrix (Placeholder)}

\IfFileExists{generated/llm_results_tables.tex}{
\input{generated/llm_results_tables.tex}
}{

\begin{table}[t]
\centering
\caption{Placeholder table for future LLM-backed evaluation.}
\label{tab:llm-matrix}
\begin{tabular}{lllllllll}
\toprule
Model & Provider & Ctrl WR & Sleep WR & Ctrl Acc & Sleep Acc & Ctrl Tok & Sleep Tok & $E_c$ \\
\midrule
\texttt{[[MODEL\_A]]} & \texttt{[[PROVIDER]]} & \texttt{[[TBD]]} & \texttt{[[TBD]]} & \texttt{[[TBD]]} & \texttt{[[TBD]]} & \texttt{[[TBD]]} & \texttt{[[TBD]]} & \texttt{[[TBD]]} \\
\texttt{[[MODEL\_B]]} & \texttt{[[PROVIDER]]} & \texttt{[[TBD]]} & \texttt{[[TBD]]} & \texttt{[[TBD]]} & \texttt{[[TBD]]} & \texttt{[[TBD]]} & \texttt{[[TBD]]} & \texttt{[[TBD]]} \\
\bottomrule
\end{tabular}
\end{table}
}

\subsection{Statistical Validation Plan (Placeholder)}

\begin{itemize}[leftmargin=1.25em]
\item \texttt{[[N\_SEEDS]]} random seeds per model/condition
\item Report mean $\pm$ std and confidence intervals
\item Paired comparison over identical task instances
\item Ablations:
  \begin{itemize}[leftmargin=1.25em]
  \item no Dream
  \item no entropy trigger (interval only)
  \item no pruning
  \item no graph summaries (rules only)
  \item raw summaries vs.\ heuristic rules
  \end{itemize}
\end{itemize}

\subsection{Results Interpretation Text (Final Paper Template)}

When replacing the placeholder tables with LLM-backed results, the final manuscript should include a concise interpretation subsection that reports both effect size and tradeoffs. A recommended structure is:

\begin{itemize}[leftmargin=1.25em]
\item \textbf{Primary effect}: State whether Sleep improves task success and by how much (absolute and relative deltas).
\item \textbf{Efficiency effect}: Report token budget ratio and $E_c$, clarifying whether gains come from compression, improved decisions, or both.
\item \textbf{Uncertainty dynamics}: Compare entropy trajectories (or token-level uncertainty if available), and note whether Sleep reduces drift late in episodes.
\item \textbf{Domain heterogeneity}: Identify domains where Sleep helps most/least and hypothesize why (e.g., recurring motifs vs.\ one-off tasks).
\item \textbf{Ablation interpretation}: Explain whether Dream, entropy triggers, or pruning contribute materially to final gains.
\item \textbf{Failure modes}: Note cases where Sleep harms performance (over-compression, stale heuristics, incorrect pruning) and quantify frequency.
\end{itemize}

Example reporting language (to be replaced with actual numbers): ``Across $N$ seeds on SWE-bench Lite, Agentic Sleep improved resolve rate by $X$ absolute points while reducing inference tokens to $Y\%$ of the control condition, yielding a consolidation efficiency of $E_c = Z$. Gains were concentrated in multi-file issues with recurring failure patterns, while one-shot localization-heavy tasks showed smaller improvements.''

\section{Discussion}

The prototype demonstrates a practical pattern: memory performance can improve by \textbf{changing representation}, not only by increasing storage. The Sleep cycle acts as an adaptive compression and reset mechanism. In the simulated setting, it reduced token budget and uncertainty while improving task outcomes.

The strongest conceptual contribution is the move from:
\begin{itemize}[leftmargin=1.25em]
\item \textbf{Narrative retention} (``what happened'') to
\item \textbf{Procedural retention} (``what tends to work under conditions $X$'').
\end{itemize}

This distinction matters for long-horizon agents where repeated task motifs appear across turns or sessions.

\section{Limitations}

This work is currently limited by the prototype setup:

\begin{itemize}[leftmargin=1.25em]
\item The decision policy is simulated and has access to a hidden ``best action'' via the environment.
\item The tasks are synthetic proxies, not real interactive game environments.
\item Entropy is computed over simulated action distributions, not true next-token distributions.
\item The semantic graph is lightweight and domain-specific; richer relation extraction is not yet implemented.
\item Dream self-play is a heuristic simulator, not a learned world model.
\end{itemize}

Therefore, current results should be treated as \textbf{architecture validation}, not definitive evidence of LLM capability gains.

\section{Conclusion}

Agentic Sleep provides a concrete framework for long-horizon agent stability through off-cycle consolidation, heuristic abstraction, and synthetic hardening. The current prototype demonstrates that such a design can improve performance-per-token in a controlled simulation. The next stage is direct LLM integration with the same memory loop and evaluation protocol, enabling a rigorous test of whether structural abstraction outperforms raw-context retention in real autonomous agent workloads.

\appendix

\section{Reproduction (Current Prototype)}

Run the prototype benchmark:

\begin{verbatim}
python3 -m agentic_sleep.cli --games 20 --turns 50
\end{verbatim}

Key files:
\begin{itemize}[leftmargin=1.25em]
\item \texttt{agentic\_sleep/engine.py}
\item \texttt{agentic\_sleep/consolidation.py}
\item \texttt{agentic\_sleep/dream.py}
\item \texttt{agentic\_sleep/memory.py}
\item \texttt{agentic\_sleep/experiment.py}
\end{itemize}

\section{Suggested Next Fill-Ins for the LLM Version}

\begin{itemize}[leftmargin=1.25em]
\item Replace simulated \texttt{thought} string with model-generated rationale summary (optional logging only)
\item Add logprob-based entropy extraction
\item Add real tool traces (e.g., coding agent actions) into \texttt{EpisodeStep.tags}
\item Persist MTEM/LTSM across sessions
\item Evaluate on code repair / planning / debugging benchmarks
\end{itemize}

\bibliographystyle{plainnat}
\bibliography{references}

\end{document}
